\documentclass[runningheads]{llncs}

\usepackage[T1]{fontenc}
\usepackage{graphicx}
\usepackage{booktabs}
\usepackage{multirow}
\usepackage{amsmath}
\usepackage{url}

\begin{document}

\title{Evaluating Evidence-Grounded Emotion Recognition in\\
Social Dialogue with Retrieval-Augmented\\
and Graph-Augmented LLMs}

\titlerunning{Evidence-Grounded Emotion Recognition in Social Dialogue}

\author{Seungwoo Kim}
\authorrunning{S. Kim}

\institute{University of Technology Sydney, Sydney, NSW 2007, Australia\\
\email{fo.swkim@gmail.com}}

\maketitle

% ─────────────────────────────────────────────────────────────────────────────
\begin{abstract}
Emotion recognition in conversation (ERC) requires models to ground their
predictions in conversational evidence rather than relying solely on
parametric knowledge. This paper presents a systematic evaluation of
six context-construction strategies for LLM-based ERC: LLM-only,
full-context, BM25 retrieval-augmented generation (RAG), dense RAG,
text-only RAG, and graph-augmented RAG. Beyond standard classification
metrics, we introduce a claim-level faithfulness evaluation that measures
whether model-generated explanations are supported by the retrieved
evidence. Experiments on MELD and EmoryNLP reveal complementary insights.
On MELD, LLM-only inference achieves the highest accuracy (0.618) but
near-zero faithfulness (0.014), indicating near-total reliance on
parametric knowledge; graph-augmented RAG achieves the highest macro
F1 (0.491) among retrieval methods while maintaining competitive
faithfulness (0.222). Ablation analysis identifies same-speaker history
as the most critical graph component ($-$0.053 macro F1 on removal).
On EmoryNLP, severe class imbalance ($>$70\% of instances in two classes)
causes all methods to collapse to predicting only the dominant class
(macro F1 $\approx$0.12), revealing a fundamental limitation of
zero-shot prompting under extreme imbalance. Together, these findings
highlight faithfulness as a necessary complement to accuracy for
evaluating LLM-based ERC systems and motivate class-balanced evaluation
protocols for imbalanced dialogue datasets.

\keywords{Emotion Recognition in Conversation \and
Retrieval-Augmented Generation \and Graph-Augmented Retrieval \and
Evidence-Grounded Reasoning \and Faithfulness Evaluation \and
Large Language Models}
\end{abstract}

% ─────────────────────────────────────────────────────────────────────────────
\section{Introduction}
\label{sec:intro}

Emotion recognition in conversation (ERC) is a central task in social
and behavioural computing, enabling applications in dialogue analysis,
mental health support, human-robot interaction, and education
technology~\cite{poria2019meld,zahiri2018emorynlp}. Unlike
sentence-level emotion classification, ERC requires the model to
integrate contextual signals spanning multiple utterances, speaker
relationships, and discourse dynamics.

Large language models (LLMs) have demonstrated remarkable capabilities
in natural language understanding, and recent studies have begun
applying them to ERC with strong results~\cite{zhao2023chatgpt,wang2023instructerc}.
However, a fundamental question remains: when an LLM predicts an
emotion, is the prediction grounded in the conversational evidence
actually provided, or does it reflect the model's internal parametric
knowledge about how dialogue typically unfolds? In socially sensitive
settings, the ability to trace a prediction back to specific
conversational evidence is as important as classification accuracy
itself.

Retrieval-augmented generation (RAG)~\cite{lewis2020rag} offers a
natural mechanism for grounding LLM outputs by explicitly supplying
relevant context at inference time. For ERC, retrieval can select the
most emotionally relevant prior utterances, reducing noise from
irrelevant context. Graph-augmented retrieval extends this idea by
representing dialogue as a structured graph of speaker and emotional
relationships, enabling more principled evidence selection that reflects
the social dynamics of conversation.

Despite growing interest in RAG for dialogue, prior work has
predominantly focused on accuracy metrics, leaving the faithfulness of
model reasoning under-explored. This paper addresses this gap by
jointly evaluating predictive performance and evidence faithfulness
across six context-construction strategies.

\noindent\textbf{Contributions.} This paper makes the following contributions:
\begin{enumerate}
    \item We present a comparative evaluation of six LLM-based ERC
    methods spanning no-retrieval, flat retrieval, and graph-augmented
    retrieval baselines on the MELD and EmoryNLP datasets.

    \item We introduce a claim-level faithfulness evaluation framework
    in which model explanations are decomposed into atomic claims and
    each claim is verified against the retrieved evidence by an
    LLM-as-judge, yielding a continuous faithfulness score.

    \item We demonstrate a systematic faithfulness gap: LLM-only and
    full-context methods achieve near-zero faithfulness ($<$0.04),
    while retrieval-augmented methods raise faithfulness by an order of
    magnitude ($>$0.21), regardless of classification accuracy.

    \item We conduct an ablation study on graph components,
    finding that same-speaker history is the most critical structural
    element, contributing a $-$0.053 macro F1 drop when removed.
\end{enumerate}

% ─────────────────────────────────────────────────────────────────────────────
\section{Related Work}
\label{sec:related}

\subsection{Emotion Recognition in Conversation}

ERC has been extensively studied using recurrent~\cite{hazarika2018conversational}
and graph-based neural architectures~\cite{ghosal2019dialoguegcn}.
DialogueGCN~\cite{ghosal2019dialoguegcn} models speaker-specific and
inter-speaker dependencies via graph convolutional networks.
Subsequent work has incorporated commonsense knowledge graphs and
multi-modal signals to improve minority-class prediction.
More recently, pre-trained language models fine-tuned on dialogue data
have established strong baselines~\cite{ghosal2020cosmic}.

\subsection{LLMs for Social Dialogue Understanding}

The advent of instruction-tuned LLMs has motivated a shift from
fine-tuning to prompting-based ERC. Zhao et al.~\cite{zhao2023chatgpt} evaluate
ChatGPT on several ERC benchmarks and find competitive zero-shot
performance. InstructERC~\cite{wang2023instructerc} further
demonstrates that instruction-tuned LLMs can match or exceed
fine-tuned baselines when provided with appropriate dialogue context.
However, these studies focus exclusively on accuracy and do not examine
whether predictions are grounded in the supplied evidence.

\subsection{Retrieval-Augmented Generation}

RAG~\cite{lewis2020rag} augments LLM inputs with passages retrieved
from an external corpus. For dialogue tasks, retrieval can be applied
within the conversation window to select contextually relevant
utterances. BM25~\cite{robertson2009bm25} is a standard lexical
retrieval baseline, while dense passage retrieval (DPR)~\cite{karpukhin2020dpr}
uses learned encoders for semantic matching.
Gao et al.~\cite{gao2023survey} provide a comprehensive survey of RAG approaches
and highlight the importance of retrieval quality for downstream
generation fidelity.

\subsection{Graph-Augmented Retrieval for Dialogue}

Graph representations of dialogue context have been explored for
capturing inter-utterance and inter-speaker dependencies beyond
sequential order. Relation-aware graph attention networks have been
applied to ERC to model the speaker dynamics that influence emotional
expression. In this work, we adopt a graph construction that encodes
same-speaker history, emotion-shift transitions, and contextual
adjacency to provide structured evidence for LLM prompting.

\subsection{Faithfulness and Evidence-Grounded Evaluation}

Faithfulness evaluation originated in summarisation research, where
Maynez et al.~\cite{maynez2020faithfulness} showed that abstractive models frequently
hallucinate facts not present in the source.
For RAG systems, Es et al.~\cite{es2023ragas} propose RAGAS, a reference-free
evaluation suite that measures faithfulness by checking whether generated
statements are entailed by the retrieved context. Our faithfulness
evaluation follows a similar claim-verification paradigm but is adapted
to the conversational emotion recognition setting, where each atomic
claim in the model's explanation is verified against the retrieved
dialogue evidence.

% ─────────────────────────────────────────────────────────────────────────────
\section{Methodology}
\label{sec:method}

\subsection{Problem Formulation}

Let $D = \{u_1, u_2, \ldots, u_T\}$ denote a dialogue, where each
utterance $u_t = (s_t, x_t)$ consists of speaker $s_t$ and text $x_t$.
Given a target utterance $u_\tau$ and its conversational context
$C_\tau = \{u_t : t < \tau\}$, the task is to predict the emotion
label $y_\tau \in \mathcal{Y}$ and produce a natural language
explanation $e_\tau$ grounded in the available evidence.

\subsection{Context-Construction Methods}

We compare six methods that differ in how they construct the evidence
provided to the LLM:

\paragraph{LLM-only.}
Only the target utterance $u_\tau$ is included in the prompt. No
contextual evidence is supplied, so any prediction must rely on the
model's parametric knowledge.

\paragraph{Full-context LLM.}
The complete preceding context $C_\tau$ is appended to the prompt,
providing maximal information but potentially introducing noise from
utterances unrelated to the target emotion.

\paragraph{BM25 RAG.}
A BM25 index~\cite{robertson2009bm25} is built over $C_\tau$, and the
top-$k$ utterances most lexically similar to the target are retrieved
as evidence.

\paragraph{Dense RAG.}
Utterances in $C_\tau$ are encoded by a sentence transformer
\texttt{all-MiniLM-L6-v2}. The top-$k$ utterances with the highest
cosine similarity to the target embedding are retrieved.

\paragraph{Text-only RAG.}
The context window utterances (the $k$ most recent utterances) are
provided as plain text without retrieval scoring or graph structure.
This baseline controls for the effect of structure versus content.

\paragraph{Graph-augmented RAG.}
A dialogue graph $G_\tau = (V, E)$ is constructed over $C_\tau$, where
each node $v \in V$ represents an utterance. Edges encode three types
of relationships: (i) \emph{same-speaker links} connecting utterances
by the same speaker to capture within-speaker emotional continuity;
(ii) \emph{emotion-shift links} connecting adjacent utterances where
a notable emotional transition is detected; and (iii) \emph{sequential
adjacency links} connecting temporally adjacent utterances.
Graph-based evidence is then serialised into natural language and
included in the prompt alongside the retrieved utterances.

\subsection{Prompting Strategy}
\label{sec:prompting}

All methods share a consistent evidence-constrained prompt format
designed to encourage faithfulness. The system instruction states that
the model must base its explanation \emph{only} on the provided
evidence and must not introduce external knowledge. The user message
contains three blocks: \texttt{[Dialogue Context]}, \texttt{[Target
Utterance]}, and \texttt{[Retrieved Evidence]}. The model is instructed
to output a JSON object with three fields:
\texttt{predicted\_emotion}, \texttt{explanation}, and
\texttt{evidence\_used}. The evidence-constrained prompt includes an
explicit rule set:

\begin{enumerate}
    \item The explanation \emph{must} directly quote or cite specific
    utterances from \texttt{[Retrieved Evidence]}.
    \item The model must not introduce background knowledge or
    assumptions not present in the evidence.
    \item If evidence is insufficient, this must be stated explicitly.
    \item The \texttt{evidence\_used} field must list the exact phrases
    relied upon.
\end{enumerate}

\subsection{Claim-Level Faithfulness Evaluation}
\label{sec:faithfulness}

Standard accuracy metrics cannot measure whether a prediction is
supported by the provided evidence. We introduce a two-step
claim-level faithfulness evaluation:

\paragraph{Step 1 — Claim decomposition.}
Given a model explanation $e_\tau$, an LLM judge decomposes it into a
set of atomic, independently verifiable claims
$\{c_1, c_2, \ldots, c_n\}$.

\paragraph{Step 2 — Claim verification.}
Each claim $c_i$ is presented to the judge together with the dialogue
context and retrieved evidence. The judge assigns one of three verdicts:
\emph{supported} (score 1.0), \emph{partially supported} (score 0.5),
or \emph{unsupported} (score 0.0).

The faithfulness score for a prediction is the mean score over all
extracted claims:
\begin{equation}
    \text{Faithfulness}(e_\tau) = \frac{1}{n}\sum_{i=1}^{n} \text{score}(c_i)
\end{equation}

Both decomposition and verification are performed by \texttt{gpt-4o-mini}
with temperature 0 for reproducibility. The supported ratio and
unsupported ratio are also reported as secondary faithfulness metrics.

% ─────────────────────────────────────────────────────────────────────────────
\section{Experimental Setup}
\label{sec:setup}

\subsection{Datasets}

\paragraph{MELD.} The Multimodal EmotionLines Dataset~\cite{poria2019meld}
is derived from the TV series \textit{Friends} and contains 13{,}708
utterances across 1{,}433 multi-party dialogues. Each utterance is
annotated with one of seven emotion labels: \textit{anger, disgust,
fear, joy, neutral, sadness,} and \textit{surprise}. The test split
contains 2{,}610 utterances across 280 dialogues.

\paragraph{EmoryNLP.} The EmoryNLP dataset~\cite{zahiri2018emorynlp}
is also sourced from \textit{Friends} but uses a different annotation
scheme with eight labels: \textit{mad, joyful, peaceful, neutral, sad,
scared, powerful,} and \textit{surprised}. The test split contains
1{,}328 utterances across 85 scenes. Due to the severe class imbalance
in EmoryNLP (neutral and joyful together account for $>$70\% of
instances), macro F1 is the primary evaluation metric.

\subsection{Model and Inference}

All experiments use \texttt{gpt-4o-mini} as the inference model with
temperature 0 and a maximum of 512 output tokens, ensuring deterministic
outputs across runs. The same model is used as the faithfulness judge.
Inference is parallelised with eight concurrent API workers. All
experiments use 500 instances sampled from the test split to control
API cost while maintaining statistical coverage.

\subsection{Retrieval Settings}

For BM25 and dense retrieval, the top $k=5$ utterances are retrieved
from the context window $C_\tau$. For dense retrieval,
\texttt{sentence-transformers/all-MiniLM-L6-v2} is used as the
sentence encoder. The context window for text-only RAG consists of
the three most recent utterances. For graph-augmented RAG, the graph
is built over the same context window, and evidence is serialised using
a natural-language template.

\subsection{Evaluation Metrics}

\paragraph{Classification metrics.} We report accuracy, macro-averaged
F1 (macro F1), weighted F1, worst-class F1, and per-class F1 for each
method. Macro F1 is the primary metric because it treats all emotion
classes equally regardless of frequency.

\paragraph{Faithfulness metrics.} We report the mean faithfulness
score, supported ratio (fraction of fully supported claims), and
unsupported ratio (fraction of unsupported claims) averaged across
all evaluated instances.

% ─────────────────────────────────────────────────────────────────────────────
\section{Results and Analysis}
\label{sec:results}

\subsection{MELD Classification Performance}

Table~\ref{tab:overall-performance} reports classification performance
on the MELD test set. Graph RAG achieves the highest macro F1 (0.491)
among all methods, outperforming Text-only RAG by 0.011 and BM25 RAG
by 0.038. LLM-only achieves the highest accuracy (0.618) but, as shown
in Section~\ref{sec:faith-results}, this does not correspond to
faithful reasoning.

\begin{table}[t]
\centering
\caption{Classification performance on MELD test set (500 instances).
Best result per column is \textbf{bold}.}
\label{tab:overall-performance}
\begin{tabular}{lcccc}
\toprule
Method & Accuracy & Macro F1 & Weighted F1 & Worst-class F1 \\
\midrule
LLM-only          & \textbf{0.618} & 0.485 & \textbf{0.621} & 0.292 \\
Full-context LLM  & 0.574          & 0.491 & 0.586          & 0.328 \\
BM25 RAG          & 0.526          & 0.453 & 0.539          & 0.298 \\
Dense RAG         & 0.529          & 0.454 & 0.541          & 0.292 \\
Text-only RAG     & 0.538          & 0.480 & 0.548          & 0.333 \\
Graph RAG         & 0.570          & \textbf{0.491} & 0.580 & 0.320 \\
\bottomrule
\end{tabular}
\end{table}

\subsection{Label-wise Performance}

Table~\ref{tab:perclass} reports per-class F1 scores on MELD.
Graph RAG shows the strongest performance on contextually dependent
classes: \textit{disgust} (0.400), \textit{neutral} (0.655), and
\textit{surprise} (0.581). These are classes for which speaker
relationship history is particularly informative, supporting the value
of graph-structured evidence. Minority classes such as \textit{fear}
and \textit{disgust} remain difficult for all methods, with F1 scores
below 0.40 across the board.

\begin{table}[t]
\centering
\caption{Per-class F1 on MELD test set. Best per class is \textbf{bold}.}
\label{tab:perclass}
\begin{tabular}{lcccccccc}
\toprule
Method & ang & dis & fear & joy & neu & sad & sur \\
\midrule
LLM-only         & 0.487 & 0.328 & 0.292 & 0.588 & \textbf{0.750} & 0.411 & 0.537 \\
Full-context LLM & \textbf{0.525} & 0.394 & \textbf{0.328} & 0.576 & 0.667 & 0.420 & 0.524 \\
BM25 RAG         & 0.518 & 0.298 & 0.308 & 0.529 & 0.598 & 0.414 & 0.504 \\
Dense RAG        & 0.531 & 0.316 & 0.292 & 0.532 & 0.600 & 0.405 & 0.503 \\
Text-only RAG    & 0.524 & 0.333 & 0.333 & 0.542 & 0.582 & \textbf{0.478} & 0.569 \\
Graph RAG        & 0.475 & \textbf{0.400} & 0.320 & \textbf{0.563} & 0.655 & 0.444 & \textbf{0.581} \\
\bottomrule
\end{tabular}
\end{table}

\subsection{Faithfulness Analysis}
\label{sec:faith-results}

Table~\ref{tab:faithfulness} reports faithfulness metrics.
The most striking finding is the near-total absence of faithfulness in
non-retrieval methods: LLM-only achieves a mean faithfulness score of
only 0.014, with 98.4\% of claims rated unsupported. Full-context LLM
improves marginally to 0.033, still with 96.3\% unsupported claims.
This reveals that even when the full dialogue is provided, the LLM
predominantly relies on its internal parametric knowledge rather than
the supplied context.

Among retrieval methods, BM25 RAG achieves the highest faithfulness
(0.269), followed by Dense RAG (0.256), Graph RAG (0.222), and
Text-only RAG (0.213). The ordering reflects the precision of retrieval:
BM25's lexical matching tends to retrieve utterances that share
vocabulary with the target, making them more directly citable. Graph RAG
outperforms Text-only RAG in faithfulness (+0.009), suggesting that
structured relational evidence provides slightly more citable content
than unselected context.

\begin{table}[t]
\centering
\caption{Faithfulness evaluation on MELD test set (500 instances).
Best result per column is \textbf{bold}.}
\label{tab:faithfulness}
\begin{tabular}{lccc}
\toprule
Method & Faithfulness$\uparrow$ & Supported$\uparrow$ & Unsupported$\downarrow$ \\
\midrule
LLM-only          & 0.014          & 0.013          & \textbf{0.984} \\
Full-context LLM  & 0.033          & 0.028          & 0.963 \\
BM25 RAG          & \textbf{0.269} & \textbf{0.199} & 0.662 \\
Dense RAG         & 0.256          & 0.180          & 0.668 \\
Text-only RAG     & 0.213          & 0.152          & 0.726 \\
Graph RAG         & 0.222          & 0.164          & 0.720 \\
\bottomrule
\end{tabular}
\end{table}

\subsection{EmoryNLP Results}
\label{sec:emorynlp-results}

Table~\ref{tab:emorynlp-performance} reports classification performance
on EmoryNLP. In stark contrast to MELD, all six methods collapse to
predicting only the dominant \textit{neutral} class, yielding
worst-class F1 of 0.0 across the board. This complete minority-class
failure is attributable to the extreme class imbalance: \textit{neutral}
accounts for approximately 35\% and \textit{joyful} for approximately
36\% of EmoryNLP instances, leaving six minority classes at fewer than
10\% combined. All methods achieve macro F1 in the range 0.105--0.124,
with differences smaller than the measurement noise. LLM-only achieves
the highest accuracy (0.977) precisely because it overwhelmingly predicts
\textit{neutral}, illustrating the deceptiveness of accuracy as a metric
under severe imbalance. Faithfulness scores for llm\_only (0.044) and
full\_context\_llm (0.061) mirror the MELD pattern: explanations for
non-retrieval methods are almost entirely unsupported by evidence.
These results confirm that the class-imbalance problem in EmoryNLP
prevents meaningful comparison of retrieval strategies under the current
zero-shot prompting setup and underline the importance of addressing
imbalance (e.g., via label-balanced sampling or class-weighted loss)
in future work.

\begin{table}[t]
\centering
\caption{Classification performance on EmoryNLP test set (500 instances).
All methods collapse to predicting \textit{neutral}; macro F1 is near
chance due to severe class imbalance.}
\label{tab:emorynlp-performance}
\begin{tabular}{lcccc}
\toprule
Method & Accuracy & Macro F1 & Weighted F1 & Faithfulness \\
\midrule
LLM-only          & \textbf{0.977} & \textbf{0.124} & \textbf{0.988} & 0.044 \\
Full-context LLM  & 0.919          & 0.120          & 0.958          & 0.061 \\
BM25 RAG          & 0.751          & 0.107          & 0.858          & --- \\
Dense RAG         & 0.723          & 0.105          & 0.839          & --- \\
Text-only RAG     & 0.746          & 0.107          & 0.854          & --- \\
Graph RAG         & 0.728          & 0.105          & 0.843          & --- \\
\bottomrule
\end{tabular}
\end{table}

\subsection{Ablation Study}
\label{sec:ablation}

Table~\ref{tab:ablation} reports the ablation study on graph components.
Removing same-speaker history causes the largest degradation in both
accuracy ($-$0.076) and macro F1 ($-$0.053), confirming that
within-speaker emotional continuity is the most informative relational
signal. In contrast, removing audio features slightly \emph{improves}
accuracy (+0.004), suggesting that audio-derived features introduce
noise in the text-only inference pipeline used here. Removing
emotion-shift links causes negligible performance change, indicating
that emotion-shift transitions contribute less than speaker identity to
the graph evidence.

\begin{table}[t]
\centering
\caption{Ablation results on MELD test set. $\Delta$ macro F1 is
relative to Graph RAG (full).}
\label{tab:ablation}
\begin{tabular}{lcccc}
\toprule
Variant & Accuracy & Macro F1 & $\Delta$ Macro F1 & Faithfulness \\
\midrule
Graph RAG (full)        & 0.570 & \textbf{0.491} & --- & 0.222 \\
$-$Same-speaker history & 0.494 & 0.438 & $-$0.053 & 0.181 \\
$-$Emotion-shift links  & 0.564 & 0.488 & $-$0.003 & 0.234 \\
$-$Audio features       & \textbf{0.574} & 0.495 & +0.004 & 0.235 \\
\midrule
Text-only RAG           & 0.538 & 0.480 & $-$0.011 & 0.213 \\
\bottomrule
\end{tabular}
\end{table}

% ─────────────────────────────────────────────────────────────────────────────
\section{Discussion}
\label{sec:discussion}

\paragraph{The faithfulness--accuracy trade-off.}
Our results reveal a clear trade-off between accuracy and faithfulness.
LLM-only achieves the highest accuracy (0.618) because it can draw on
extensive parametric knowledge about how social emotions manifest in
television dialogue, a genre well-represented in LLM training corpora.
However, its 98.4\% unsupported claim rate indicates that this accuracy
is not grounded in the conversational evidence. Retrieval-augmented
methods constrain the model to reason from specific utterances,
slightly reducing accuracy but raising faithfulness by an order of
magnitude.

\paragraph{Why accuracy decreases with retrieval.}
Retrieval introduces a precision--recall trade-off in evidence selection.
When the retrieved context does not contain the most relevant emotional
cue (e.g., a distant utterance that established the conversational
dynamic), the model may be misled by irrelevant evidence. This explains
why BM25 RAG and Dense RAG underperform text-only methods in accuracy
despite achieving higher faithfulness. Graph RAG partially mitigates
this by using relational structure to prioritise same-speaker history,
achieving the best macro F1 among retrieval methods.

\paragraph{Same-speaker history as the dominant signal.}
The ablation study identifies same-speaker history as the most
informative graph component. This is consistent with psychological
findings that an individual's emotional trajectory is auto-correlated:
a speaker who was angry in a recent turn is more likely to remain
angry or transition to related states~\cite{hazarika2018conversational}.
Structural encoding of this continuity allows the graph to surface
more relevant evidence than unstructured flat retrieval.

\paragraph{EmoryNLP class-imbalance collapse.}
EmoryNLP exposes a fundamentally different challenge. With \textit{neutral}
and \textit{joyful} together comprising over 70\% of instances, the
constrained zero-shot prompting paradigm fails to overcome the prior
imbalance: every method effectively predicts only \textit{neutral},
achieving macro F1 near $1/8 = 0.125$ (the random baseline for eight
classes if one class dominates). This is \emph{not} a failure of graph
augmentation specifically — all methods, including LLM-only and
full-context, exhibit identical behaviour. The result points to a
methodological prerequisite: class-imbalanced ERC datasets require
label-balanced sampling or class-weighted prompting before retrieval
strategy comparisons are meaningful.

\paragraph{Limitations.}
Several limitations should be noted. First, the faithfulness evaluator
is itself an LLM, introducing potential evaluation bias. Second,
all methods use the same evidence-constrained prompt, which may
inherently limit faithfulness for methods with less citable evidence
(e.g., LLM-only). Third, the graph construction relies on heuristic
edge weights that may not capture all relevant dialogue dynamics.
Fourth, EmoryNLP results are uninformative under zero-shot prompting
due to class imbalance; future work should address this with balanced
sampling. Finally, the current experiments are limited to 500 instances
per dataset; full-dataset evaluation may reveal different patterns for
rare emotion classes.

% ─────────────────────────────────────────────────────────────────────────────
\section{Conclusion}
\label{sec:conclusion}

This paper presented a comparative evaluation of retrieval-augmented
and graph-augmented LLM methods for evidence-grounded emotion recognition
in social dialogue. A key contribution is the introduction of a
claim-level faithfulness evaluation framework that measures whether
model explanations are supported by retrieved evidence.

Our experiments on MELD demonstrate that non-retrieval LLM methods,
despite achieving high accuracy, produce explanations that are almost
entirely unsupported by the conversational evidence (98.4\% unsupported
claims). Graph-augmented RAG achieves the best macro F1 among retrieval
methods and improves faithfulness over flat text retrieval, with
same-speaker history identified as the most critical structural
component. On EmoryNLP, results reveal a distinct but equally important
failure mode: under extreme class imbalance ($>$70\% of instances in
two classes), all zero-shot methods collapse to predicting the dominant
class, yielding macro F1 near chance regardless of retrieval strategy.
This finding motivates class-balanced evaluation and sampling protocols
as a prerequisite for meaningful comparison of evidence-grounded methods.

These findings highlight faithfulness as a necessary complement to
accuracy for evaluating LLM-based ERC systems in socially sensitive
settings. Future work should investigate learned graph construction
strategies, label-balanced prompting for imbalanced datasets,
cross-dataset generalisation, and the relationship between faithfulness
and human-perceived explanation quality.

% ─────────────────────────────────────────────────────────────────────────────
\bibliographystyle{splncs04}
\bibliography{references}

\end{document}
